%%
%% Berliner Hochschule für Technik --  Abschlussarbeit
%%
%% Kapitel 3
%%
%%

\chapter{Evaluation geeigneter KI-Modelle}
\section{Technische Grundlagen der KI-gestützten Reizerkennung}

Die App kann zwei Arten von Reizen verwenden. Erstens können vordefinierte Bilder eingesetzt werden, deren Kategorie und Label bereits bekannt sind. Zweitens können Nutzende eigene Bilder oder Kamerabilder verwenden. Dann muss die App erkennen, ob und welche rauchbezogenen Objekte oder Szenen enthalten sind. Diese zweite Variante erfordert ein Klassifikationsmodell und ist technisch deutlich anspruchsvoller.

Der Laborstudie kann ein für die App nutzbares Datenmodell entnommen werden. Ein Zielreiz wird mit einer kleinen Menge möglicher Antworten kombiniert. Für diese Arbeit lässt sich das als Relation zwischen Bild-ID, Reizkategorie, Antworttyp, Labeltext, Distraktor und Bedingung modellieren. Ein Eintrag kann beispielsweise die Kategorie \enquote{Rauch}, das Label \enquote{Rauch}, einen Distraktor wie \enquote{Packung} und die Bedingung \enquote{Cue-Labeling} enthalten. Für Cue-Matching wird dagegen ein weiterer Reiz derselben oder einer kontrolliert ähnlichen Kategorie angeboten. Eine solche Struktur ist technisch einfach. Wenn die gleichen Bilder, Kategorien und Antwortoptionen in beiden Bedingungen verwendet werden, kann später geprüft werden, ob Bedingungsunterschiede tatsächlich auf Labeling und nicht auf unterschiedlich stark auslösende Bilder zurückgehen.

Statistical Learning unterscheidet Klassifikationsprobleme von Regressionsproblemen. Bei der Klassifikation wird eine Kategorie vorhergesagt, etwa ob ein Bild eine Zigarette, Rauch, ein Feuerzeug oder keinen relevanten Reiz enthält \cite{james_islr_2021}. Für CueLens ist nicht allein die Genauigkeit entscheidend. Ein falsch positiver Treffer kann dazu führen, dass einer Person ein Rauchreiz oder Label angeboten wird, obwohl das Bild keinen relevanten Reiz enthält. Das kann irritieren und die Akzeptanz senken. Ein falsch negativer Treffer verhindert dagegen eine mögliche Intervention in einer relevanten Situation. Daraus folgt, dass Sensitivität, Spezifität, Präzision und Recall getrennt betrachtet werden müssen. Die F1-Kennzahl kann als harmonisches Mittel von Präzision und Recall nützlich sein, darf aber nicht die inhaltliche Bewertung der Fehlertypen ersetzen. Formal können die Bildentscheidungen in einer Konfusionsmatrix dargestellt werden. Für eine Reizklasse gilt dann: Präzision ist \(TP/(TP+FP)\), Recall beziehungsweise Sensitivität ist \(TP/(TP+FN)\), und \(F1=2\cdot(\mathrm{Precision}\cdot\mathrm{Recall})/(\mathrm{Precision}+\mathrm{Recall})\). Bei ungleichen Klassenhäufigkeiten ist zusätzlich zu prüfen, ob ein hoher Gesamtwert vor allem durch die häufige Klasse \enquote{kein Rauchreiz} entsteht.

Information-Retrieval-Grundlagen sind für die Labelauswahl relevant. Wenn ein Bild mehrere mögliche Reize enthält, muss die App passende Beschreibungen oder Vergleichsreize auswählen. Manning et al. beschreiben hierfür allgemeine Prinzipien wie Tokenisierung, Indexierung, Gewichtung und Rangordnung von Treffern \cite{manning_ir_2008}. Übertragen auf CueLens bedeutet dies, dass Labels und Vergleichsbilder nicht beliebig präsentiert werden, sondern aus einer strukturierten Reiz- und Labelmenge stammen sollten. Eine einfache Datenstruktur kann aus Reizkategorien, Synonymen, Bild-IDs, Labeltexten und Schwellenwerten bestehen. Für jedes erkannte Objekt wird ein Score gespeichert. Die App kann anschließend nur Labels anbieten, deren Score einen Mindestwert überschreitet. Dadurch wird die Auswahl nachvollziehbar, testbar und später austauschbar.

On-Device-Klassifikation begrenzt die Modellgröße, den Energieverbrauch und die Rechenzeit. Für die technische Evaluation sind daher neben Klassifikationsmetriken auch Speicherbedarf, Inferenzzeit, Modellgröße, Energieverbrauch und Android-Kompatibilität zu dokumentieren. Ein Modell, das auf einem Entwicklungsrechner hohe Genauigkeit erreicht, ist für diese Arbeit ungeeignet, wenn es auf realen Smartphones zu langsam, zu groß oder zu instabil ist.

\newpage
\null
\newpage
\null
\newpage
\null
\newpage
\null
